% Q2: Parallel RLC Step Response — Current source with R, L, C in parallel
\documentclass[border=10pt]{standalone}
\usepackage[american voltages, american currents]{circuitikz}
\usepackage{amsmath}
\renewcommand{\familydefault}{\sfdefault}

\begin{document}
\begin{circuitikz}[
    line width=0.8pt,
    every node/.style={font=\sffamily}
]

% Current source (left, vertical, arrow pointing up)
\draw (0,0) to[I, l=$I_s$] (0,4);

% Top rail
\draw (0,4) -- (2,4);

% Junction dots
\fill (2,4) circle (2pt);
\fill (5,4) circle (2pt);
\fill (8,4) circle (2pt);

% Branch 1: Resistor
\draw (2,4) to[R, l_=$R$] (2,0);

% Branch 2: Inductor with current label
\draw (5,4) to[L, l_=$L$, i>_=$i_L$] (5,0);

% Branch 3: Capacitor
\draw (8,4) to[C, l_=$C$] (8,0);

% Top rail connections
\draw (2,4) -- (5,4) -- (8,4);

% Bottom rail
\draw (0,0) -- (2,0) -- (5,0) -- (8,0);

% Bottom rail dots
\fill (2,0) circle (2pt);
\fill (5,0) circle (2pt);

% Voltage polarity v(t) — gap between source and R
\draw (1,4) to[open, v^=$v(t)$] (1,0);

\end{circuitikz}
\end{document}
